\chapter{2004年重庆卷（理）}
\section{单选题}
\begin{problem}
函数 \( y = \sqrt{\log_3(3x - 2)} \) 的定义域是（\quad）

\choices
    {\(\{1, +\infty\}\)}
    {\(\left(\frac{2}{3}, +\infty\right)\)}
    {\(\left[\frac{2}{3}, 1\right]\)}
    {\(\left(\frac{2}{3}, 1\right]\)}
\end{problem}

\begin{problem}
设复数 \(z = 1 + \sqrt{2}\i\)，则 \(Z^2 - 2Z =\)（\quad）

\choices
    {\(-3\)}
    {3}
    {-3i}
    {3i}
\end{problem}

\begin{problem}
圆 \(x^{2} + y^{2} - 2x + 4y + 3 = 0\) 的圆心到直线 \(x - y = 1\) 的距离为（\quad）

\choices
    {2}
    {\(\frac{\sqrt{2}}{2}\)}
    {1}
    {\(\sqrt{2}\)}
\end{problem}

\begin{problem}
不等式 \(x + \frac{2}{x + 1} > 2\) 的解集是（\quad）

\choices
    {\((-1,0)\cup (1, + \infty)\)}
    {\((-\infty, -1) \cup (0, 1)\)}
    {\((-1,0)\cup (0,1)\)}
    {\((- \infty, -1) \cup (1, +\infty)\)}
\end{problem}

\begin{problem}
\(\sin 163^{\circ}\sin 223^{\circ} + \sin 253^{\circ}\sin 313^{\circ} =\)（\quad）

\choices
    {\(-\frac{1}{2}\)}
    {\( \frac{1}{2} \)}
    {\(-\frac{\sqrt{3}}{2}\)}
    {\(\frac{\sqrt{3}}{2}\)}
\end{problem}

\begin{problem}
若向量 \(\bs{a}\) 与 \(\bs{b}\) 的夹角为 \(60^{\circ}\)，\(\left|\bs{b}\right| = 4\)，\(\left(\bs{a} + 2\bs{b}\right) \cdot \left(\bs{a} - 3\bs{b}\right) = -72\)，则向量 \(\bs{a}\) 的模为（\quad）

\choices
    {2}
    {4}
    {6}
    {12}
\end{problem}

\begin{problem}
一元二次方程 \(ax^2 + 2x + 1 = 0 (a \neq 0)\) 有一个正根和一个负根的先分不必要条件是（\quad）

\choices
    {\(a < 0\)}
    {\(a > 0\)}
    {\(a < -1\)}
    {\(a > 1\)}
\end{problem}

\begin{problem}
设 P 是 \( 60^{\circ} \) 的二面角 \( \alpha - l - \beta \) 内一点，\( PA \perp \) 平面 \( \alpha \) ，\( PB \perp \) 平面 \( \beta \) ，A, B 为垂点，PA = 4, PB = 2，则 AB 的长为（\quad）

\choices
    {\(2\sqrt{3}\)}
    {\(2 \sqrt{3}\)}
    {\(2\sqrt{7}\)}
    {\(4\sqrt{2}\)}
\end{problem}

\begin{problem}
若数列 \( \{n_{n}\} \) 是等差数列，首项 \( a_{1}>0, a_{2000}+a_{2001}>0, a_{2003}\cdot a_{2004}<0 \) ，则使前 n 项和 \( S_{n}>0 \) 成立的最大自然数 n 是（\quad）

\choices
    {4005}
    {4006}
    {4007}
    {4008}
\end{problem}

\begin{problem}
已知双曲线 \(\frac{x^2}{a^2} - \frac{y^2}{b^2} = 1 (a > 0, b > 0)\) 的左、右焦点分别为 \(F_1, F_2\)，点 \(P\) 在双曲线的右支上，且 \(|PF_1| = 4|PF_2|\)，则此双曲线的离心率 \(e\) 的最大值为（\quad）

\choices
    {\( \frac{4}{3} \)}
    {\( \frac{5}{3} \)}
    {2}
    {\( \frac{7}{3} \)}
\end{problem}

\begin{problem}
某校高三年级举行一次演讲赛共有 10 位同学参赛，其中一班有 3 位，二班有 2 位，其它班有 5 位. 若采用抽签的方式确定他们的演讲顺序，则一班有 3 位同学恰好被排在一起 (指演讲序号相连). 而二班的 2 位同学没有被排在一起的概率为（\quad）

\choices
    {\(\frac{1}{10}\)}
    {\(\frac{1}{20}\)}
    {\(\frac{1}{40}\)}
    {\(\frac{1}{120}\)}
\end{problem}

\begin{problem}
设一汽车在前进途中要经过4个路口，汽车在每个路口遇到绿灯的概率为 \(\frac{3}{5}\)，遇到红灯 (禁止通行) 的概率为 \(\frac{1}{4}\). 假定汽车只在遇到红灯或到达目的地才停止前进，表示停车时已经通过的路口数，求: (1) (哪棵树的分布列及期望 \(E\)); (2) 停车时最多已通过 3 个路口的概率.
\end{problem}

\section{填空题}
\begin{problem}
若在 \((1 + ax)^5\) 的展开式中 \(x^3\) 的系数为 \(-80\)，则 \(a = \)  \fillinblank{}.
\end{problem}

\begin{problem}
曲线 \( y = 2 - \frac{1}{2} x^2 \) 与 \( y = \frac{1}{4} x^3 - 2 \) 在交点处切线的夹角是 \fillinblank{}. (用弧度数作答)
\end{problem}

\begin{problem}
如图 \( P_{1} \) 是一块半径为 1 的半圆形纸板，在 \( P_{1} \) 的左下端剪去一个半径为 \( \frac{1}{2} \) 的半圆后得到图形 \( P_{2} \) . 然后依次剪去一个更小半圆（其直边有一个被剪掉半圆的半径）得圆形 \( P_{3} \) ，\( P_{4} \) ，\( \cdots \) ，\( P_{n} \) ，\( \cdots \) ，记纸板 \( P_{n} \) 的面积为 \( S_{n} \) ，则 \( \lim_{n\to\infty}S_{n}= \)  \fillinblank{}.
\end{problem}

\begin{problem}
对任意实数 \(k\)，直线 \(y = kx + b\) 与椭圆 \(\left\{ \begin{array}{l} x = \sqrt{3} + 2\cos \theta, \\ y = 1 + 4\sin \theta, \end{array} \right.\) (\(0 \leqslant \theta \leqslant 2\pi\)) 恰有一个公共点，则 \(b\) 取值范围是 \fillinblank{} \fillinblank{}.
\end{problem}

\section{解答题}
\begin{problem}
求函数 \( y = \sin^4 x + 2\sqrt{3}\sin x\cos x - \cos^4 x \) 的取小正周期和取小值；并写出该函数在 \([0, \pi]\) 上的单调递增区间.
\end{problem}

\begin{problem}
设一汽车在前进途中要经过 4 个路口，汽车在每个路口遇到绿灯的概率为 \(\frac{3}{5}\)，遇到红灯 (禁止通行) 的概率为 \(\frac{1}{4}\). 假定汽车只在遇到红灯或到达目的地才停止前进，\(\xi\) 表示停车时已经通过的路口数，求: (1) \(\xi\) 的概率的分布列及期望 \(EG\); (2) 停车时最多已通过 3 个路口的概率.
\end{problem}

\begin{problem}
如图，四棱锥 \(P - ABCD\) 的底面是正方形，\(PA \perp\) 底面 \(ABCD\)，\(AE \perp PD\)，\(EF // CD\)，\(AM = EF\). (1) 证明: \(MF\) 是异轴直线 \(AB\) 与 \(PC\) 的公垂线; (2) 若 \(PA = 3AB\)，求直线 \(AC\) 与平面 \(EAM\) 所成角的正弦值.
\end{problem}

\begin{problem}
设函数 \( f(x) = x(x - 1)(x - a) (a > 1) \). (1) 求导数 \( f'(x) \) ，并证明 \( f(x) \) 有两个不同的极值点 \( x_{1}, x_{2} \) ; (2) 若不等式 \( f(x_{1}) + f(x_{2}) \leqslant 0 \) 成立，求 \( \alpha \) 的取值范围.
\end{problem}

\begin{problem}
设 \( p > 0 \) 是一常数，过点 \( Q(2p,0) \) 的直线与抛物线 \( y^{2} = 2px \) 交于相异两点 \( A, B \)，以线段 \( AB \) 为直径作圆 \( H(H \) 为圆心). 试证抛物线顶点在圆 \( H \) 的圆周上; 并求圆 \( H \) 的面积最小时直线 \( AB \) 的方程.
\end{problem}

\begin{problem}
设数列 \(\{a_{n}\}\) 满足 \(a_1 = 2, a_{n+1} = a_n + \frac{1}{a_n}\) (\(n = 1, 2, 3, \ldots\)). (1) 证明 \(a_{n} > \sqrt{2n + 1}\) 对一切正整数 \(n\) 成立； (2) 令 \( b_{n} = \frac{a_{n}}{\sqrt[n]{n}} (n = 1, 2, \cdots) \)，判断 \( b_{n} \) 与 \( b_{n+1} \) 的大小，并说明理由.
\end{problem}

